\documentclass[10pt,a4paper]{amsart}
\usepackage[margin=19mm]{geometry}
\usepackage{amsmath,amssymb,mathtools}
\usepackage[scr=rsfs,cal=euler]{mathalfa}
\usepackage{xfrac}
\usepackage[unicode,hidelinks,bookmarks=false]{hyperref}
\newcommand{\ie}{i.e.\ }
\newcommand{\cf}{cf.\ }
\newcommand{\resp}{respectively\ }
\newcommand{\Subsection}{\subsection}
\providecommand{\nobreakdash}{\nobreak\hbox{-}}
\setlength{\emergencystretch}{2em}
\multlinegap0pt
\usepackage{mathtools}
\usepackage[all]{xy}
\usepackage{amssymb}
\usepackage[shortlabels]{enumitem}
\usepackage{tikz}
\usepackage{xcolor}
\newtheorem{axiom}{Axiom}
\newtheorem{lemma}{Lemma}
\newtheorem{theorem}{Theorem}
\newcommand{\C}{\mathcal C}
\newcommand{\Comp}{\operatorname{Comp}}
\newcommand{\Fill}{\mathcal F}
\newcommand{\id}{\operatorname{id}}
\newcommand{\one}{1}
\title{\textbf{Noetherian forms of free non-symmetric operads}}
\author{Zurab Janelidze}
\date{Source: arXiv:2609.10501v1 (CC BY 4.0)}
\begin{document}
\maketitle

\noindent\emph{Source and licence.} \textbf{Noetherian forms of free non-symmetric operads}. Version arXiv:2609.10501v1 (CC BY 4.0), distributed by arXiv under the Creative Commons Attribution 4.0 International licence, http://creativecommons.org/licenses/by/4.0/, as recorded in the arXiv licence metadata for this exact version and shown on the abstract page https://arxiv.org/abs/2609.10501v1. Full licence notice: \texttt{licenses/arxiv-2609.10501-CC-BY-4.0.txt}.\par\smallskip

\noindent\textit{Editorial scope.} Counted unit: the article's theorem concl:symmetric-representation (source0.tex lines 4407--4415). The article's complete original proof of it is delivered with it (source0.tex lines 4416--4554, one \texttt{proof} environment of 139 lines). No mathematics is written, repaired or re-derived: the statement and the proof are the article's own lines, in the article's own notation, and no retained line is substituted or re-flowed.  Retained uncounted as the source-local context the proof needs: lines 400--404; lines 412--435; lines 448--458; lines 1034--1042; lines 1133--1141; lines 1857--1862; lines 4376--4380; lines 4394--4397.  Nothing was omitted from any retained span, so the package carries no omission record.  No retained line contains a citation, so the wrapper carries no bibliography and no bibliography entry.  No retained line contains a figure environment, an image-inclusion command or drawing code, so no image is delivered and the package declares no asset dependency.  Editorial formatting only, in addition: an AMS \texttt{amsart} preamble that restates the article's own declarations and macros used by the retained lines, namely the environments \texttt{axiom}, \texttt{lemma}, \texttt{theorem} on separate counters, together with the standard packages those lines need.  The retained environments print in this document as Axiom 1, Lemma 1, Theorem 1 in order of appearance, whereas the article prints the same items with the numbering of its own full document.  The complete native source of the article as distributed in the arXiv e-print for this exact version is bundled unchanged as \texttt{sources/209924/source0.tex}.
\begin{axiom}[pushout stability]\label{ax:pointpushout}
The class of epimorphisms whose pullbacks along all but at most one
point are isomorphisms is stable under pushout along strong
monomorphisms.
\end{axiom}

\begin{axiom}[terminal pullback factorizations]\label{ax:U}
For every object $A$ and point $j:\one\to B$, the following
category, denoted by $\Fill(A,j)$, has a terminal object. An
object of $\Fill(A,j)$ is a triple $(C,d,u)$, where
$u:A\to C$ is a strong monomorphism, $d:C\to B$ is an
epimorphism, and the square
\begin{equation}\label{eq:fillsquare}
\vcenter{\xymatrix@C=4em@R=2.5em{
A\ar[r]^u\ar[d]_{q_A}&C\ar[d]^d\\
 \one\ar[r]_j&B
}}
\end{equation}
is a pullback. A morphism from $(C',d',u')$ to $(C,d,u)$ is
a morphism $h:C'\to C$ making the following diagram commute
\[
\xymatrix@C=3.5em@R=2.5em{
&&C'\ar[dl]_h\ar@/^1pc/[ddl]^{d'}\\
A\ar[r]^u\ar[d]_{q_A}\ar@/^1pc/[urr]^{u'}&C\ar[d]^d&\\
\one\ar[r]_j&B&
}
\]
Thus $hu'=u$ and $dh=d'$.
Composition and identities are those of $\C$.
\end{axiom}

\begin{axiom}[diagonals]\label{ax:O}
Every commutative square of strong monomorphisms
\[
\xymatrix@C=4em@R=2.5em{
W \ar[r]^{a} \ar[d]_{b} & X \ar[d]^{c} \\
Y \ar[r]_{d} & Z
}
\]
admits either a diagonal $h:X\to Y$ satisfying $ha=b$ and $dh=c$,
or a diagonal $k:Y\to X$ satisfying $kb=a$ and $ck=d$.
\end{axiom}

\begin{lemma}[Factorization of constant composites]\label{lem:constantfactor}
Let $u:A\to X$ be a strong monomorphism, and let $j:\one\to X/u$
be the point in the pushout defining $\pi_u$.
For every point $k\ne j$ of $X/u$, the pullback of $\pi_u$ along
$k$ is an isomorphism.
If $w:W\to X$ is strong monic with nonterminal domain and $\pi_u w$
is constant, then $\pi_u w=jq_W$ and there is a unique $t:W\to A$
such that $ut=w$.
\end{lemma}

\begin{lemma}[One-component pushouts]\label{lem:componentpushout}
Push out an indecomposable strong monomorphism $i:A\to X$ along any epimorphism
$a:A\to B$, writing $c:X\to Y$ and $j:B\to Y$ for the new arrows.
Then $Y\ne\one$, $j$ is an indecomposable strong monomorphism, and the component
decomposition of $Y$ consists exactly of $j$ and the strong monomorphisms
$ck$ for $k\in\Comp(X)\setminus\{i\}$. These strong monomorphisms are
distinct. The restriction of $c$ at $i$ is $a$, and its restriction
at every other component is an identity.
\end{lemma}

\begin{theorem}[Finite filling from singleton filling]\label{thm:finite}
Every finite marked cocone of constant morphisms has a unique universal filling. It replaces exactly
the marked empty leaves by their assigned objects. It may be constructed
by successive singleton fillings in any order, and the result, including
its structure arrows, is independent of that order.
\end{theorem}

\begin{equation}\label{concl:corolla-symmetry}
 \operatorname{Aut}_{\C}(\phi)
 \xrightarrow{\ \sim\ }\operatorname{Sym}(\Comp(\phi)),
 \qquad g\longmapsto\bigl([u]\longmapsto[gu]\bigr).
\end{equation}

\begin{equation}\label{concl:symmetric-matching}
 (K,p,\theta),\qquad
 \theta:X/K\xrightarrow{\sim}Y|_p,
\end{equation}

\begin{theorem}\label{concl:symmetric-representation}
An essentially small, locally small category satisfies the modified
axioms above if and only if it is equivalent to
$\mathcal T_L^{\mathrm{sym}}$ for a ranked set $L$. The labels are
recovered as the isomorphism classes of corollas, and their arities
are the numbers of components. Under the equivalence, epimorphisms
are prunings followed by tree isomorphisms, and strong monomorphisms
are entire-subtree occurrences preceded by tree isomorphisms.
\end{theorem}

\begin{proof}
We first describe the reconstruction, retaining comparison
isomorphisms at the steps where the original proof gave equalities.
Terminal maps are epic: apply Axiom~\ref{ax:U} with
$j=\id_1$, so that the upper map of its pullback square is
invertible and its right map is epic. Points are strong
monic because they are split monic. It follows also that a
strong monomorphism into $1$, or an epimorphism out of $1$,
is an isomorphism.

For each $X$, its finite strong-subobject order is a tree branching
order: it has greatest element $[\id_X]$, and the subobjects above
any fixed subobject form a chain by Axiom~\ref{ax:O}.
Postcomposition by a strong monomorphism $A\to X$ identifies
the subobjects of $A$ with those below its class. Thus each position
has a unique finite path of component inclusions to the root.

Push out $q_A$ along a strong monomorphism $u:A\to X$, obtaining
the elementary collapse $X\to X/u$. Its fibre at the new point is
$A$, and Axiom~\ref{ax:pointpushout} makes every other point fibre
terminal. Consequently, a strong monomorphism with nonterminal
domain which becomes constant under this collapse factors through
$u$. This is the argument of Lemma~\ref{lem:constantfactor}, with
terminal domains understood up to isomorphism. Factoring the image
of a disjoint subobject, and applying point detection to its epic
factor, then shows that this subobject remains strong monic.

Every noninvertible epimorphism has a nonterminal point fibre.
Collapsing that fibre factors the epimorphism through a
noninvertible elementary collapse. Repetition strictly increases
the quotient class of the original domain, so finiteness makes it
terminate. Every epimorphism is therefore a sequence of elementary
collapses followed by an isomorphism.

The component calculation in
Lemma~\ref{lem:componentpushout} applies to subobject classes:
inverse image preserves properness, and direct image after inverse
image is the identity. A pushout at one component consequently
replaces that component and preserves the other components. Applying
this successively shows that an epimorphism with nonterminal target
induces a bijection of components and has epic component
restrictions. Collapsing every component gives a corolla $r(X)$ and
an epimorphism $\rho_X:X\to r(X)$. The pushout property identifies
root corollas along epimorphisms with nonterminal targets.

Induction on the finite strong-subobject tree now classifies
quotients as effective cuts. A quotient with nonterminal target
is obtained by its component restrictions; a quotient with
terminal target collapses the whole tree. Collapses at disjoint
positions commute by their pushout property, and a collapse at
an ancestor absorbs collapses below it. Thus the successive
collapses may be recorded by pairwise incomparable source
positions with nonterminal domains. These positions are
recovered from the quotient as the maximal source positions
with nonterminal domains whose images have terminal domains.
This proves uniqueness as well as existence of the cut.

We next explain the role of filling. A terminal singleton filling
has terminal fibres at its unmarked points: collapsing any additional
nonterminal fibres gives an epic comparison $e:C\to C'$ to
another filling. Its comparison $h:C'\to C$ back to the
terminal filling satisfies $he=\id_C$ by terminality.
Epicity of $e$ then gives $eh=\id_{C'}$. The cut description
excludes a noninvertible collapse with this property. The unmarked
points therefore lift uniquely, so singleton filling can be iterated
at finitely many distinct points. The proof of
Theorem~\ref{thm:finite} gives a terminal simultaneous filling.
Its comparisons are epic: a proper strong image would lie in a
component and would remain proper after the epic map to the
nonterminal base.

In particular, prescribe an object at every component point of a
corolla. Every complete filling is isomorphic to the terminal one.
Indeed, its comparison has isomorphic restrictions at all components,
so every point pullback is an isomorphism and the modified detection
axiom applies. For a nullary corolla this argument uses detection
with an empty set of points. Thus complete fillings are uniquely
isomorphic when the corolla and the component inclusions are fixed.

Suppose $X$ and $Y$ have isomorphic root corollas. Condition
\eqref{concl:corolla-symmetry} assigns a unique root-corolla
isomorphism to each component bijection. The universal property of
complete fillings then gives
\begin{equation}\label{concl:recursive-isomorphisms}
 \operatorname{Iso}_{\C}(X,Y)\cong
 \coprod_{\sigma:\Comp(X)\xrightarrow{\sim}\Comp(Y)}
 \prod_{i\in\Comp(X)}\operatorname{Iso}_{\C}(X_i,Y_{\sigma(i)}).
\end{equation}
For clarity, prescribed restrictions $a_i:X_i\to Y_{\sigma(i)}$
extend to the unique comparison $h$ in
\[
\xymatrix@C=3.2em@R=2.4em{
X_i\ar[r]^{u_i}\ar[d]_{a_i}&X\ar[r]^{\rho_X}\ar[d]^h
 &r(X)\ar[d]^\beta\\
Y_{\sigma(i)}\ar[r]_{v_{\sigma(i)}}&Y\ar[r]_{\rho_Y}&r(Y).
}
\]
Both rows specify complete fillings over the same corolla after
transport along $\beta$, so $h$ is an isomorphism. If the root
corollas differ, there are no isomorphisms. Recursion now recovers
exactly the isomorphisms of labeled trees with arbitrary child
permutations.

Finally, normal factorization, with representatives chosen for cuts
and subobjects, gives
\begin{equation}\label{concl:symmetric-hom}
 \C(X,Y)\cong
 \coprod_{K\text{ a cut of }X}\ \coprod_{p\in
 \operatorname{Sub}_{\mathrm{s}}(Y)}
 \operatorname{Iso}_{\C}(X/K,Y_p).
\end{equation}
These are exactly the triples~\eqref{concl:symmetric-matching}.
Direct image on positions preserves composition by uniqueness of
strong images. A matching's positional function determines its
target position, its cut, and its surviving tree isomorphism, so
\eqref{concl:symmetric-hom} respects composition. Complete fillings
construct every finite tree, which proves essential surjectivity.

Conversely, the triples define a category because their positional
functions compose as just described and distinguish morphisms.
Prunings are epic, being surjective on positions. A proper
occurrence is not epic, since its collapse and the corresponding
constant map agree on that subtree and differ at the ambient root.
Factorization therefore identifies the epimorphisms as prunings
followed by isomorphisms. Occurrences are strong monic: in a lifting
square, surjectivity of the epic map on positions forces the other
map's image into the specified subtree, where it restricts uniquely.
These observations identify the two intrinsic arrow classes.
Restricting a pruning gives the required
pullback, and performing it inside an ambient subtree gives the
bicartesian pushout. The surviving isomorphisms restrict in the
pullback and extend uniquely in the pushout. Filling a marked empty
leaf has the required terminal property: a comparison erases the
other cut fibres, and its surviving isomorphism is fixed on the
inserted tree and on its context. The remaining axioms follow from
finite cuts and positions, nesting of overlapping subtrees, detection
of a nonempty cut by its point fibres, and the full permutation group
of a corolla.
\end{proof}

\end{document}
